% ===== PRÉAMBULE CANONIQUE — à coller VERBATIM en tête de CHAQUE .tex =====
\documentclass[11pt,a4paper]{article}
\usepackage[utf8]{inputenc}\usepackage[T1]{fontenc}\usepackage{lmodern}
\usepackage[french]{babel}
\usepackage{amsmath,amssymb,mathtools}\usepackage{icomma}
\usepackage[version=4]{mhchem}          % \ce{...} : équations & formules
\usepackage{siunitx}\sisetup{locale=FR} % \qty{}{}, \si{}, \ang{}
\usepackage{chemfig}                    % \chemfig{...} : structures développées
\usepackage{xcolor}\usepackage{colortbl}
\usepackage{tikz}\usepackage{pgfplots}\pgfplotsset{compat=1.18}
\usepackage[most]{tcolorbox}\usepackage{enumitem}\usepackage{array,booktabs}
\usepackage[a4paper,margin=2.2cm]{geometry}
\usetikzlibrary{arrows.meta,calc,angles,quotes,positioning,patterns,backgrounds,shapes.geometric}
\definecolor{primary}{RGB}{222,49,99}\definecolor{primarydark}{RGB}{150,22,55}
\definecolor{defcol}{RGB}{0,114,178}\definecolor{methcol}{RGB}{52,92,148}
\definecolor{excol}{RGB}{0,135,90}\definecolor{attncol}{RGB}{200,40,55}
\tcbset{boxbase/.style={enhanced,breakable,boxrule=0.9pt,arc=2.5pt,
  left=3mm,right=3mm,top=2.3mm,bottom=2.3mm,fonttitle=\bfseries,coltitle=white}}
% --- boîtes TUTORIEL (noms EXACTS, ne pas renommer) ---
\newtcolorbox{cle}[1][]{boxbase,colback=defcol!6,colframe=defcol,
  title={\textsc{À retenir}\ifx\relax#1\relax\relax\else\textmd{ : \itshape #1}\fi}}
\newtcolorbox{methode}[1][]{boxbase,colback=methcol!7,colframe=methcol,sharp corners,
  title={\textsc{Méthode}\ifx\relax#1\relax\relax\else\textmd{ --- \itshape #1}\fi}}
\newtcolorbox{exemplebox}[1][]{boxbase,colback=excol!5,colframe=excol,
  title={\textsc{Exemple}\ifx\relax#1\relax\relax\else\textmd{ : \itshape #1}\fi}}
\newtcolorbox{piege}[1][]{boxbase,colback=attncol!6,colframe=attncol,
  title={\textsc{Piège}\ifx\relax#1\relax\relax\else\textmd{ : \itshape #1}\fi}}
\newtcolorbox{remarque}[1][]{boxbase,colback=black!4,colframe=black!45,
  title={\textsc{Remarque}\ifx\relax#1\relax\relax\else\textmd{ : \itshape #1}\fi}}
% --- boîtes EXERCICE / CORRIGÉ (noms EXACTS, auto-comptées) ---
\newtcolorbox[auto counter]{exo}[1][]{boxbase,colback=primary!4,colframe=primary,
  title={\textsc{Exercice}~\thetcbcounter\ifx\relax#1\relax\relax\else\textmd{ --- \itshape #1}\fi}}
\newtcolorbox[auto counter]{corrige}[1][]{boxbase,colback=excol!4,colframe=excol,
  title={\textsc{Corrigé}~\thetcbcounter\ifx\relax#1\relax\relax\else\textmd{ --- \itshape #1}\fi}}
\newcommand{\R}{\mathbb{R}}\newcommand{\N}{\mathbb{N}}\newcommand{\Z}{\mathbb{Z}}\newcommand{\ds}{\displaystyle}
% (un \usepackage{fancyhdr}+en-tête est possible ; il est ignoré côté web)
% =========================================================================
\begin{document}

\begin{exo}[métal à valence variable et anion simple]
La valence du métal est imposée par le chiffre romain. Croise-la avec celle de l'anion.
\begin{enumerate}[label=\alph*)]
  \item fer(III) \ce{Fe^3+} et \ce{Cl-}
  \item cuivre(II) \ce{Cu^2+} et \ce{O^2-}
  \item étain(IV) \ce{Sn^4+} et \ce{O^2-}
\end{enumerate}
\end{exo}

\begin{exo}[métal et anion polyatomique]
Croise les valences, puis parenthèse l'anion polyatomique s'il apparaît au moins deux fois.
\begin{enumerate}[label=\alph*)]
  \item \ce{Ca^2+} et \ce{CO3^2-}
  \item \ce{Al^3+} et \ce{SO4^2-}
  \item strontium \ce{Sr^2+} et iodate \ce{IO3-}
\end{enumerate}
\end{exo}

\begin{exo}[sels d'ammonium]
L'ammonium \ce{NH4+} est un cation polyatomique de valence 1. Croise, puis parenthèse \ce{NH4} s'il apparaît
au moins deux fois.
\begin{enumerate}[label=\alph*)]
  \item \ce{NH4+} et sulfite \ce{SO3^2-}
  \item \ce{NH4+} et phosphate \ce{PO4^3-}
  \item \ce{NH4+} et carbonate \ce{CO3^2-}
\end{enumerate}
\end{exo}

\begin{exo}[du nom complet à la formule]
Identifie le cation et l'anion (avec leur valence) d'après le nom, puis écris la formule brute.
\begin{enumerate}[label=\alph*)]
  \item sulfate d'aluminium
  \item nitrate de calcium
  \item phosphate de sodium
\end{enumerate}
\end{exo}

\begin{exo}[distinguer des anions voisins]
Une lettre change l'anion, donc la formule. Écris chaque formule brute.
\begin{enumerate}[label=\alph*)]
  \item sulfite de magnésium
  \item sulfate de magnésium
  \item thiosulfate de sodium
\end{enumerate}
\end{exo}

\end{document}
